\begin{textsample}{POS Dim 2 – am – Score 22.00 – am/am\_inf\_4.txt}  \label{ex:f2_pos_159}
Percurso metodológico da construção do referencial curricular amazonense para a Educação Infantil A \textbf{trajetória} metodológica pela busca da construção do Referencial Curricular Amazonense para a Educação Infantil, \textbf{alinhada} à Base Nacional Comum Curricular, iniciou seguida à indicação paritária de 122 articuladores no interior e 14 na capital, contando ainda com o acordo de parceria \textbf{técnica} firmado com o \textbf{regime} de colaboração entre Secretaria de Estado de Educação e Qualidade do Ensino – Seduc/AM e União dos \textbf{Dirigentes} \textbf{Municipais} de Educação – Undime/AM.
A comissão foi composta por \textbf{profissionais} da Seduc/AM e Secretarias \textbf{Municipais} de Educação, em articulação com representantes da sociedade civil organizada, Universidade Federal do Amazonas e Universidade Estadual do Amazonas.
A interlocução com os \textbf{municípios} ocorreu por meio de encontros presenciais e a distância, mediado por fóruns, questionários, formulários de contribuições e gravação de vídeos.
Desta feita, a tônica do percurso de construção do Referencial Curricular Amazonense para a Educação Infantil se deu a partir, pesquisas com dados fornecidos pelos \textbf{municípios} sobre o trabalho desenvolvido com bebês, crianças bem pequenas e crianças pequenas.
O processo de interlocução entre os diferentes sujeitos que atuam na Educação Infantil nas instituições de educação dos \textbf{municípios} foi importante para situar o ponto de partida e para o estabelecimento de um diálogo com a organização do \textbf{currículo} previsto pela BNCC.
No dia 16 de março de 2018, aconteceu o Dia D de Mobilização da BNCC, um momento em que professores, pedagogos, \textbf{funcionários} e comunidade se reuniram para estudar o documento a fim de compreender a organização e a proposta de implementação.
A Comissão Estadual de Implementação da BNCC iniciou os trabalhos em prol da elaboração do Referencial Curricular Amazonense para a Educação Infantil com reunião de apresentação dos integrantes da comissão, organização do plano de trabalho junto às equipes, discussão sobre o termo de compromisso dos bolsistas, apresentação do cronograma geral e definição de ações futuras, que ocorreu nos dias 13, 21 e 22/3/2018.
Ainda neste mês, a equipe amazonense participou do I Encontro \textbf{Formativo} do Programa de Apoio à implementação da BNCC, em Brasília, nos dias 26, 27 e 28 de março de 2018.
No período de março a julho, a equipe redatora trabalhou na construção da versão \textbf{preliminar} deste Referencial Curricular.
Houve debates e momentos integrativos que se coadunaram para dar corpo à elaboração do documento, \textbf{alinhado} às orientações pedagógicas firmadas com a equipe da BNCC.
No mês de abril, os trabalhos continuaram priorizando os ajustes para \textbf{efetivação} de espaços para a acomodação da equipe de redatores.
Seguiu -se a escrita dos campos de experiências, evidenciando os direitos de aprendizagens nos objetivos sinalizados na BNCC e nos acrescidos pela equipe redatora, contemplando a regionalidade e considerando as diferentes faixas etárias.
De 17 a 30 de maio de 2018 foi disponibilizado um formulário online de contribuições ao documento por meio das ações dos articuladores.
Houve participação de 83 \% dos \textbf{municípios} do estado, gerando 5.766 contribuições analisadas pelos redatores.
Também houve a gravação de um vídeo explicativo no Centro de Mídias da Seduc/AM sobre o trabalho em desenvolvimento na Educação Infantil, esclarecendo o processo em construção do \textbf{currículo}, além das concepções que regem a etapa, disponibilizado na internet para acesso público.
As diferentes interlocuções com setores da Seduc/AM, SEMED, Conselho Estadual de Educação do Amazonas – CEE/AM, Universidade Estadual do Amazonas – UEA e Universidade Federal do Amazonas – Ufam, além da participação na web conferência com os temas “ Organização curricular na Educação Infantil ” e “ Transição da Educação Infantil para o Ensino Fundamental ”, com equipe do MEC, acrescentaram à discussão importantes olhares.
Contamos também, no mês de maio de 2018, com a apreciação do professor dr.
Roberto Sanches Mubarac Sobrinho, da UEA, o que possibilitou o acréscimo de tópicos que tornaram o trabalho mais integrador em uma perspectiva de continuidade das discussões necessárias para a escrita de um documento que \textbf{atenda} as especificidades dos \textbf{municípios} amazonenses.
Em junho de 2018, houve vários momentos que contribuíram significativamente para as ações dos redatores, como a participação da equipe amazonense no II Encontro \textbf{Formativo} do Programa de Apoio à implementação da Base Nacional Comum Curricular – BNCC, em Brasília, nos dias 12 e 13 de junho de 2018, em Brasília/DF, com \textbf{carga} \textbf{horária} de 16 horas.
Ressaltamos que após esse encontro, a equipe redatora amazonense retornou com outros subsídios no que compete à prática da escrita, objetivando produzir um documento claro, coeso e acessível.
O \textbf{currículo} em construção contou com o apoio dos conselheiros e assessores pedagógicos do Conselho Municipal de Educação de Manaus – CME/Manaus.
Também aconteceu a web conferência “ \textbf{Regime} de colaboração e a construção do \textbf{currículo} ” com a equipe do MEC.
De 12 a 20 de julho de 2018, os \textbf{municípios} realizaram a Semana de Estudos e Reflexão para que toda a \textbf{rede} pudesse conhecer a versão \textbf{preliminar} do Referencial Curricular Amazonense.
Assim, os educadores apropriam -se do material e se prepararam para realizar contribuições na fase de consulta pública.
Os trabalhos ficaram ainda mais intensos com encaminhamentos das contribuições dos \textbf{municípios}, que foram analisados e respondidos pela equipe redatora.
Nesse mesmo mês aconteceu também a web conferência “ \textbf{Currículo} na Educação Infantil ” com a equipe do MEC, gerando momentos de reflexões e aprimoramento na escrita.
No dia 2 de agosto de 2018 ocorreu o lançamento da versão \textbf{preliminar} do documento, com a adesão de 100 \% dos \textbf{municípios} à construção coletiva.
Houve também o primeiro Encontro Estadual dos Articuladores, no qual receberam \textbf{capacitação} e orientações sobre o \textbf{currículo} e a consulta pública.
No período de 6 de agosto a 6 de setembro, o documento em versão \textbf{preliminar} esteve disponível na plataforma online do Ministério da Educação para a apreciação e opinião pública.
Este foi um momento onde todas as pessoas puderam exercer sua cidadania democrática, contribuindo com a educação do estado.
No dia 2 de agosto também ocorreu a web conferência “ Direitos de Aprendizagem ” com a equipe do MEC.
Ocorreu em agosto, na Ufam, um encontro de orientação sobre os caminhos necessários para a construção de um documento que \textbf{atenda} as especificidades dos bebês e crianças das instituições amazonense, com a professora doutora Michelle Bissoli.
Em setembro de 2018, o texto foi submetido aos leitores críticos convidados : professora doutora Michelle de Freitas Bissoli – Universidade Federal do Amazonas, professora mestre Eliseanne Lima da Silva – Instituto Federal do Amazonas.
Também houve a contribuição da professora doutora Arlene Araújo Nogueira – Universidade Federal do Amazonas, com a leitura crítica do campo de experiências, Espaços, tempos, quantidades, relações e transformações.
Os meses de setembro e outubro foram \textbf{destinados} ao trabalho de sistematização e finalização do documento.
Em 23 de outubro de 2018, a equipe redatora recebeu as contribuições do Conselho Estadual de Educação do Amazonas.
Além dessas, a equipe também analisou as contribuições dos leitores críticos e das consultas públicas, objetivando a construção de um referencial curricular que dialogue com os \textbf{municípios}, apontando os caminhos necessários para uma Educação Infantil comprometida com a aprendizagem, o desenvolvimento, a humanização e o respeito pela infância.
Em novembro de 2018, o Referencial Curricular Amazonense para a Educação Infantil foi entregue para o Conselho Municipal de Educação, responsável pela análise da etapa.
Em dezembro de 2018, o Conselho entregou a avaliação do documento para a coordenação da Educação Infantil com as considerações e encaminhamentos necessários antes de sua aprovação.
Ressalta -se, ainda, a participação de toda a equipe no III Encontro \textbf{Formativo} do Programa de Apoio à implementação da Base Nacional Comum Curricular – BNCC, em Brasília, entre os dias 12, 13 e 14 de dezembro de 2018 para debater a formação continuada de professores, \textbf{alinhada} aos novos \textbf{currículos} de referência.

% matched lemmas: alinhar, atender, capacitação, carga, currículo, destinar, dirigente, efetivação, formativo, funcionário, horário, municipal, município, preliminar, profissional, rede, regime, trajetória, técnico
\end{textsample}
